\section{Numerical results and ablations}
\label{sec:results}

All timings in this section are single-process measurements on the local
macOS 15.7.3 arm64 machine and are descriptive rather than hardware-normalized
complexity claims.

\subsection{The certified result for m=4}

The final prior has nonnegative levels
\[
 (0.5984699814893856,\;2.1264117002421385,\;4)
\]
and total level masses
\[
 (0.36699366016033996,\;0.3994285510645993,\;
   0.23357778877506074).
\]
Each positive level is split equally between its positive and negative atom.
The independent 192-bit verification used tail cutoff $A=8$, 2048 Simpson
panels per point risk, 512 curvature cells, and 567 branch nodes.  It took
49.68 seconds in the recorded clean run.

\begin{table}[ht]
\centering
\caption{Outward-rounded $m=4$ certificate.  Displayed endpoints are rounded
outward from the longer Arb output.}
\label{tab:m4}
\small
\begin{tabular}{ll}
\toprule
Quantity & Independently recomputed enclosure \\
\midrule
Bayes risk & $[0.8150887565050,\,0.8150887565548]$ \\
Worst-case risk & $[0.8150887565301,\,0.8150887615302]$ \\
Minimax risk & $[0.8150887565050,\,0.8150887615302]$ \\
Certified gap upper bound & $5.02517E-9$ \\
Branch nodes & 567 \\
Arithmetic & python-flint Arb balls, 192 bits \\
\bottomrule
\end{tabular}

\end{table}

The verifier's Bayes-risk width is $4.98\times10^{-11}$.  The branch envelope
uses a $5.0\times10^{-9}$ work width.  The certified minimax-gap upper bound is
$5.02517\times10^{-9}$, below $10^{-8}$.

\subsection{Fast deterministic sweep}

Table~\ref{tab:sweep} reports float64 discovery results.  They are not
certificates except where a separate certificate is stated.

\begin{table}[ht]
\centering
\caption{Adaptive-support float64 sweep; no ex ante parameter grid.}
\label{tab:sweep}
\small
\begin{tabular}{rrrrrr}
\toprule
$m$ & levels & atoms & Bayes risk & float gap & seconds \\
\midrule
0.25 & 1 & 2 & 0.0588158765 & 0.00e+00 & 0.003 \\
0.5 & 1 & 2 & 0.1989864336 & 0.00e+00 & 0.000 \\
1 & 1 & 2 & 0.4495995092 & 0.00e+00 & 0.001 \\
1.5 & 2 & 3 & 0.5561720123 & 3.78e-10 & 0.020 \\
2 & 2 & 3 & 0.6448008729 & 3.47e-09 & 0.030 \\
3 & 3 & 5 & 0.7506254421 & 3.62e-09 & 0.076 \\
4 & 3 & 6 & 0.8150887565 & 2.64e-09 & 0.187 \\
5 & 4 & 7 & 0.8573009664 & 3.71e-09 & 0.218 \\
8 & 6 & 11 & 0.9231345345 & 6.26e-09 & 5.149 \\
\bottomrule
\end{tabular}

\end{table}

\subsection{Dependency ablation}

Table~\ref{tab:ablation} holds the prior and endpoint-risk evaluator fixed.
The two generic global bounds exhaust the 5001-node budget.  Natural interval,
one-noise affine, Cartesian posterior-flow, and common-mean flow bounds
terminate.  These rows are float64 diagnostics; the Arb calculation supplies
the certified result.

\begin{table}[ht]
\centering
\caption{$m=4$, target separator width $10^{-8}$.}
\label{tab:ablation}
\small
\begin{tabular}{lrrrr}
\toprule
Separator & nodes & exhausted & width & seconds \\
\midrule
global Lipschitz & 5001 & yes & 2.34e-01 & 0.115 \\
global semiconvex & 5001 & yes & 5.09e-04 & 0.113 \\
natural interval & 291 & no & 8.58e-09 & 2.248 \\
one-noise affine & 427 & no & 9.95e-09 & 2.208 \\
Cartesian flow & 281 & no & 5.78e-09 & 3.170 \\
common-mean flow & 277 & no & 4.57e-09 & 0.076 \\
\bottomrule
\end{tabular}

\end{table}

The common-mean formula reduces the Cartesian-flow tree from 281 to 277 nodes.
Its much shorter runtime in this implementation comes primarily from
vectorized engineering.  Across cell
widths from approximately $3.87$ to $10^{-5}$, the ratio of common-mean to
Cartesian curvature width ranged from $0.864$ to $0.996$.  The measured
mathematical tightening is therefore modest.

\subsection{Numerical checks}

A seeded set of 100 valid symmetric priors and parameter intervals produced
no dense-sample violation of the float64 curvature enclosure. Such checks
provide falsification evidence; outward-rounded inclusion supplies the proof.
